\section{讨论}

现有结果最直接的结论是：$\beta=0.5$ 截面不能简单分为负侧普通有隙
Haldane-like 区域和正侧无隙区域。较低能量的负侧延拓态同样具有代数体关联，
但它还保留慢代数弦关联、局域边缘响应和特殊纠缠低能簇。这种组合提示临界体态
与非平庸边缘结构可能共存。

当前证据仍不足以确定负侧的最终普适类。首先，OBC 下
$E(S^z=1)-E(S^z=0)$ 容易受到边缘自由度劈裂污染，不能单独解释为体能隙。
其次，现有弦关联尺寸只到 $L=80$，虽然幂律明显优于指数模型，但还不能排除
更长尺度上的交叉行为。最后，边缘磁化和 Schmidt 低能簇是拓扑解释的重要
支持，却不能替代虚空间投影对称表示。

负侧约 $\alpha=-0.3$ 的波矢变化也需要谨慎处理。它可能对应五站点单胞下的
动量折叠和 Lifshitz 型重分配，但目前没有伴随已测普适量、对称性或拓扑不变量
的突变。类似地，正侧 $\alpha\approx0.17$ 和 $0.30$ 的响应峰反映不同
观测量对有限尺寸交叉尺度的敏感性，而不是自动构成额外相。

后续工作的优先级是：第一，使用更大的 $L$ 和 $t_{\max}$ 验证负侧低能
动力学模式；第二，将普通关联和弦关联扩展到 $L=120,160$；第三，计算 MPS
虚空间投影对称表示；第四，增加其他 $\beta$ 截面后再讨论完整二维相图。

\todo{服务器结果完成后补充 DSF 尺寸和时间标度。}
\todo{补充与无隙 SPT、拓扑 Luttinger liquid 等工作的文献比较。}
